\stagechapter{I}{Two-coder comparison}

\noindent\textbf{What this stage covers.} Whether two independent coders applying the frozen codebook to the same units agree, and how that agreement is measured, interpreted and reported. The draft thesis plan has the author code all material and a blind validation set of ten apps, a second coder \emph{sought} for applicability, minimisation and mismatch decisions (otherwise a supervisor sample), recoding of two apps after at least two weeks to estimate intra-coder consistency, per-code Cohen's $\kappa$ with counts, prevalence and Gwet's AC1, and a statement that model--human agreement is not independent human-coder reliability. The whole procedure is decision D8 (open).

\section{Why and when reliability is assessed}
\pa{mcdonald2019} separates reliability (independent application) from agreement (reconciliation by discussion) \pc{mcdonald2019}{5}{subtle differences between agreement (where two or more coders reconcile differences through discussion)} and notes that calculating IRR produces a claim about the consensus process, not a result \pc{mcdonald2019}{4}{Calculating IRR does not generate data used in results}. Its guidance: the more interpretation a predetermined code requires, the more IRR matters \pc{mcdonald2019}{16}{As a general rule with multiple coders, the more interpretation that is required to code the data}, and findings that hinge on frequency counts strengthen the case \pc{mcdonald2019}{16}{If findings hinge on frequency counts, there is a strong argument that IRR should be used}; grounded-theory style coding does not need it \pc{mcdonald2019}{15}{Grounded theory rarely, if ever, requires IRR}. Thresholds should be stated in advance, even pre-registered \pc{mcdonald2019}{17}{Papers could suggest, or even preregister, a target agreement value}. \pa{oconnor2020} gives the decision list (coders, share of data, independence, segmentation, number of codes, statistic, thresholds, reporting) and stresses that disagreements are informative \pc{oconnor2020}{4}{the content of inter- coder disagreements can be equally, if not more valuable than the ultimate degree of consistency}. \pa{braun2021} rejects reliability statistics for reflexive thematic analysis; the thesis uses the coding-reliability family, so it needs them.

\section{The coefficients}
\begin{itemize}
\item \emph{Percent agreement.} Easy but not chance-corrected; \pa{hayes2007} calls it ``flawed in nearly all important respects''. \pa{hallgren2012} shows why: coders guessing 50/50 agree half the time and coders guessing 10/90 agree 82\,\% by chance alone \pc{hallgren2012}{3}{if coders were to randomly rate 50\% of subjects as “depressed” and 50\% as “not depressed” without regard to the subject’s actual characteristics}.
\item \emph{Cohen's $\kappa$ and its relatives.} \pa{cohen1960} defines $\kappa$ for two judges and unordered categories with each judge's marginals in the chance term, and $\kappa_{\max}$ as a diagnostic of fuzzy boundaries. \pa{hallgren2012} gives variants for more than two coders (average of pairwise $\kappa$, Light; Fleiss) and warns that ``no single kappa variant corrects for both bias and prevalence'' \pc{hallgren2012}{5}{No single kappa variant corrects for both bias and prevalence}; \pa{mcdonald2019} found $\kappa$ the most common statistic in HCI and CSCW (17 of 32 papers) \pc{mcdonald2019}{10}{Cohen’s kappa was the most commonly used measure of IRR (17 out of 32 papers, 53.1\%)}.
\item \emph{Krippendorff's $\alpha$.} \pa{krippendorff2004} states conditions for a coefficient to index data reliability (interchangeable coders, units as the basis) and compares seven coefficients; \pa{hayes2007} argues that $\alpha$ meets all criteria, handles missing data and several metrics, and has a bootstrapped confidence interval \pc{hayes2007}{7}{The bootstrap sampling distribution of alpha is generated by taking a random sample of 239 pairs of judgments}; \pa{artstein2008} gives one notation for $S$, $\pi$, $\kappa$ and $\alpha$ and states that the choice follows the intended interpretation, not the observed marginals.
\item \emph{ICC.} For ordinal to ratio ratings (graded scores), \pa{hallgren2012} describes the four design factors that select an ICC variant. App-rating studies use it: \pa{stoyanov2015} (pilot ICC .57, revised), \pa{schultchen2021}, \pa{messner2020} and \pa{lin2025} (pre-set thresholds, calibration, very high ICCs).
\item \emph{AC1.} \pa{gwet2008} analyses the paradoxes of $\kappa$ (high agreement, low $\kappa$ when a category is rare or omnipresent) and proposes AC1 with a calibrated chance term and variance estimators that do not assume independent raters.
\end{itemize}

\section{The kappa debate and rare codes}
\pa{feinstein1990} shows algebraically and with tables that $\kappa$ can be low despite high observed agreement and higher for worse-looking marginals; the authors prefer to drop chance correction and design balanced challenge sets. \pa{krippendorff2004} argues the opposite: if the data lack variation, a low coefficient is correct, and high percent agreement can hide zero joint detections of a rare category. \pa{gwet2008} offers AC1; \pa{krippendorff2004} and \pa{krippendorff2019} dispute its interpretation, so AC1 is best presented as a sensitivity statistic. Privacy codes are often skewed (a dark pattern appears in few apps), so the cards agree on practice: publish the $k\times k$ tables and marginals, raw agreement, joint-presence counts and positive and negative specific agreement beside any coefficient, fix the primary statistic in advance, and design the pilot so that each code occurs often enough.

\section{Thresholds and sample sizes}
\begin{itemize}
\item \emph{Thresholds.} Krippendorff's $\alpha\ge.800$ as customary and $\ge.667$ as the lowest limit for tentative conclusions, as reported in \pa{mcdonald2019} and \pa{hallgren2012} (\pc{hallgren2012}{5}{Krippendorff (1980) provides a more conservative interpretation suggesting that conclusions should be discounted for variables with values less than 0.67}); \pa{hruschka2004}: $\ge.7$--$.8$ for rough prevalence estimates and $\ge.9$ for classification of individual cases, and cut-offs of $.8$--$.9$ in its CDC studies; \pa{mayring2000}: $\kappa>.7$; \pa{lombard2002}: conventions from mass communication; \pa{artstein2008}: $.8$ and $.7$ as computational-linguistics conventions. \pa{mchugh2012} proposes a stricter scale, but it is the author's own and its relabelling of Landis and Koch's bands is idiosyncratic, so do not cite it as a field standard.
\item \emph{Sizes.} \pa{lombard2002}: a pilot of at least 30 units and a reliability sample of at least 50 units or 10\,\% of the corpus; \pa{oconnor2020}: 10--25\,\% of units is typical \pc{oconnor2020}{5}{Depending on the size of the data set, 10–25\% of data units would be typical}; \pa{mchugh2012}: at least 30 comparisons.
\item \emph{Per code.} Pooling hides poor codes; \pa{oconnor2020} asks for the range and distribution \pc{oconnor2020}{9}{Second, pooling all codes’ reliability figures means codes with poor reliability can be “hidden”}.
\end{itemize}

\section{How the corpus studies handled it}
\begin{table}[t]
\centering\small
\caption{Reliability designs in applied studies (as described in the cards).}
\label{tab:I-practice}
\begin{tabularx}{\textwidth}{@{}P{2.6cm}P{3.2cm}YY@{}}
\toprule
\textbf{Study} & \textbf{Coders, data} & \textbf{Statistic} & \textbf{Remark}\\
\midrule
\pa{malki2024} & team coding, then first author alone & none (``did not calculate IRR'') & iterative team coding; not a frozen two-coder design\\
\pa{huckvale2019} & two coders, third adjudicator & raw agreement and $\kappa$ \emph{before} reconciliation & good reporting model\\
\pa{habib2019} & double-coded half sample & $\kappa=0.82$ & fixed template, piloted in rounds\\
\pa{zou2019} & 3 coders, 12\,\% subset & Fleiss $\kappa=.75$ & before splitting the corpus\\
\pa{wilson2016} & 3 annotators per policy & Fleiss $\kappa$ per category & segment-level presence; consolidation rule\\
\pa{oloughlin2019} & 2 raters per stage & 86.7\,\% and 89.5\,\% screening agreement & first 23 apps only for reliability\\
\pa{sunyaev2015} & two raters & Janson--Olsson $\iota=0.94$ for multiple judges & scope ratings nearly unanimous\\
\pa{hutton2018} & second rater, random subset & $\kappa$, then consensus & minimal precedent\\
\pa{lin2025} & two calibrated raters & ICC, before and after calibration & calibration on 10 apps\\
\pa{cao2024} & two coders & ``Cohen's $\kappa$ of 87.5\,\%'' & likely percent agreement under a $\kappa$ label\\
\pa{doan2025} & one coder with consultation & none & flagged as a reliability gap\\
\bottomrule
\end{tabularx}
\end{table}
\pa{campbell2013} covers the common case in which one knowledgeable coder codes everything: unitise first, code a sample blind with a second coder, then report reliability (independent) and agreement (negotiated) separately; \pa{oconnor2020} recommends a clean file with segments but without codes for the second coder \pc{oconnor2020}{7}{Once the first coder has saved the coded transcript, they can then create a second document where the data segments are visible}. \pa{nielsen1990} adds the independence principle: evaluate first, merge later.

\begin{notebox}[Disputes the thesis should not hide]
(1) \pa{lombard2002} and \pa{krippendorff2004} disagree about percent agreement, overlapping coding and whether $\pi$ and $\kappa$ are liberal or conservative; (2) the \pa{feinstein1990}/\pa{gwet2008} remedies are contested by Krippendorff; (3) \pa{mchugh2012}'s scale is its author's proposal; (4) repeated discussion rounds may only produce interpretive convergence (\pa{hruschka2004}, \pa{oconnor2020}); (5) \pa{cao2024} reports $\kappa$ as a percentage, a reminder to state the statistic, the unit and the categories.
\end{notebox}

\begin{protocolbox}[Reporting checklist for the two-coder stage]
\begin{itemize}[leftmargin=1.3em]
\item Before coding: coders and their qualifications, subset size and selection (random or stratified, seed), unit of analysis, statistic per code type (nominal, ordinal), thresholds and the rule if a code fails (\pa{lombard2002}, \pa{oconnor2020}, \pa{hallgren2012}).
\item Independence: second coder receives the clean file; no conferral before coding; training hours and calibration items recorded.
\item Results: per-code values with confidence intervals, the tables and marginals, raw agreement, joint presence, $\kappa$, $\alpha$ and AC1 for skewed codes; values \emph{before} reconciliation; resolution rule and adjudicator (\pa{huckvale2019}).
\item Interpretation: scale fixed in advance; limits (few coders, interpretive convergence); effect on later analyses.
\end{itemize}
\end{protocolbox}

\begin{thesisbox}[Implications for this thesis (stage I)]
With one person as the main coder, the human reliability claim rests on a second coder for a sample; if that is not possible, say so and treat all counts as single-coder results (\pa{mcdonald2019}, \pa{campbell2013}). Fix in the Decision Log, before the pilot, the subset share, the primary statistic and threshold per code type, and the AC1 sensitivity analysis for rare codes (the draft already names $\kappa$ and AC1). Report intra-coder consistency separately from inter-coder reliability and never describe model--human agreement as human-coder reliability.
\end{thesisbox}
